\textbf{Notation.}
In this paper, the number $C$ generically denotes a positive constant independent of the parameter $n$, the index of summation $\ell$ and the variable $z$. Moreover, the generalized hypergeometric function $_pF_q$ is defined by (see, for example, \cite[p.~404]{NIST-Handbook})
	\begin{equation}\label{pFq definition}
		{}_pF_q\bigg[\begin{matrix}
			a_1,\dots,a_p\\
			b_1,\dots,b_q
		\end{matrix};z\bigg]
		\equiv
		{}_pF_q[
		a_1,\dots,a_p;
		b_1,\dots,b_q;z]
		:=\sum_{n=0}^{\infty}\frac{(a_1)_n\cdots(a_p)_n}{(b_1)_n\cdots(b_q)_n}\frac{z^n}{n!},
	\end{equation}
	where $a_1,\dots,a_p\in\CC$ and $b_1,\dots,b_q\in\CC\setminus\ZZ_{\leqslant 0}$.

\section{Preliminary lemmas}\label{Sect. 2}
In this section, we deduce three lemmas which will be used in the sequel. The first is a sharp bound for the ratio of Pochhammer symbols.
\begin{Lemma}\label{Lemma: ratio of Pochhammer symbol}
	If $a\in\CC$ and $b\in\CC\setminus\ZZ_{\leqslant 0}$, then
	\[\biggl|\frac{(a)_n}{(b)_n}\biggr|\leqslant Cn^{\operatorname{Re}(a-b)},\qquad n\in\ZZ_{>0}.\]
\end{Lemma}
\begin{proof}
	The proof for $a\in\ZZ_{\leqslant 0}$ is trivial. If $a\in\CC\setminus\ZZ_{\leqslant 0}$, Stirling's formula implies that
		\[\lim_{n\to\infty}n^{b-a}\frac{\Gamma(a+n)}{\Gamma(b+n)}=1,\]
		which concludes that $n^{b-a}\frac{(a)_n}{(b)_n}$ is bounded uniformly for $n\in\ZZ_{>0}$.
\end{proof}

The second is a simple estimate of the function $\Phi_a(x)$, which is the ``horizontal'' generating function of Stirling numbers of real order (see \cite[Section~8]{Butzer_Kilbas_Trujillo-2003}). More properties of $\Phi_a(x)$, containing the asymptotic behavior for fixed $x\in\RR$ and $a\to\pm\infty$, are studied in \cite[Section~3.2]{VanGorder-2009}.
\begin{Lemma}\label{Phi_a(x) estimate}
	Define
	\[\Phi_a(x):=\sum_{k=1}^{\infty}\frac{k^a}{k!}x^k,\qquad a,x\in\RR.\]
	Then $\Phi_a(x)\sim x^a\me^x$, $x\to +\infty$. Thus
	\[\Phi_a(x)\leqslant Kx^a\me^x,\qquad x\geqslant 1,\]
	where $K>0$ is a constant independent of $x$.
\end{Lemma}
\begin{proof}
	Take $a_n=\frac{n^a}{n!}$ and $b_n=\frac{1}{n!}$ in \cite[p.~12, Problem~72]{Polya_Szego-vol.2}.
\end{proof}

The third is a global estimate for the confluent hypergeometric function $_1F_1$.

\begin{Lemma}\label{lem-1F1}
	Let $a,c\in\CC$. Choose $N\in\ZZ_{>0}$ such that $\operatorname{Re}(c+N+1)>0$ and define
	\begin{equation}\label{G_l definition}
			G_{\ell}(z):={}_1F_1\bigg[\begin{matrix}
				a+\ell\\
				c+\ell
			\end{matrix};z\bigg],\qquad \ell\in\ZZ_{\geqslant 0}.
	\end{equation}
	Then for $\ell\geqslant N+1$,
 \begin{equation}\label{G_l inequality}
			|G_\ell(z)|\leqslant C\me^{\gamma|z|},
	\end{equation}
	where $\gamma>0$ is a constant independent of $\ell$, $N$ and $z$.
\end{Lemma}

\begin{proof}
	Since $G_{\ell}(0)=1$, we can assume that $z\ne 0$. Recall the inequality \cite[equation~(2.3)]{Joshi_Arya-1982}
	\begin{equation}\label{1F1 inequality-1}
		\biggl|{}_1F_1\bigg[\begin{matrix}
			a_0\\
			b_0
		\end{matrix};z\bigg]\biggr|\leqslant \cos\frac{\theta}{2}\cdot
		{}_1F_1\biggl[\begin{matrix}
			|a_0|\\
			|b_0|
		\end{matrix};|z|\sec\frac{\theta}{2}\bigg]
	\end{equation}
	with $\theta=\arg(b_0)\in (-\pi,\pi)$ and the inequality \cite[p. 37, equation~(3.5)]{Carlson-1966}
	\begin{equation}\label{1F1 inequality-2}
		\me^{\frac{a_0}{b_0}z}<{}_1F_1\biggl[\begin{matrix}
			a_0\\
			b_0
		\end{matrix};z\biggr]<1-\frac{a_0}{b_0}+\frac{a_0}{b_0}\me^z,\qquad b_0>a_0>0,\quad z\ne 0.
	\end{equation}
	But when $a_0\geqslant b_0>0$ and $z>0$, we have $\frac{(a_0)_k}{(b_0)_k}\leqslant \big(\frac{a_0}{b_0}\big)^k$ since $\frac{a_0+j}{b_0+j}$ decreases with respect to~${j\geqslant 0}$. Thus
	\begin{equation}\label{1F1 inequality-3}
		{}_1F_1\biggl[\begin{matrix}
			a_0\\
			b_0
		\end{matrix};z\biggr]\leqslant\me^{\frac{a_0}{b_0}z},\qquad a_0\geqslant b_0>0,\quad z>0.
	\end{equation}
	Recall $\operatorname{Re}(c+N+1)>0$ and note that $\gamma(\ell):=\max\bigl\{1,\frac{|a+\ell|}{|c+\ell|}\bigr\}$ is bounded uniformly for $\ell\in\ZZ_{\geqslant 0}$. Therefore, a combination of the inequalities \eqref{1F1 inequality-1}--\eqref{1F1 inequality-3} claims that for $\ell\geqslant N+1$,
	\[|G_\ell(z)|\leqslant 2\gamma(\ell)\me^{\sqrt{2}\gamma(\ell)|z|}\leqslant C\me^{\gamma|z|},\]
	where
	\[\gamma:=\sup_{\ell\in\ZZ_{\geqslant 0}}\sqrt{2}\gamma(\ell)\geqslant \sqrt{2}.\]
	This completes the proof.
\end{proof}

\begin{Remark}
	Lemma \ref{lem-1F1} can be easily generalized to the following form. Let $a_1,\dots,a_p,b_1,\dots$, $b_p\in\CC$, and let $N\in\ZZ_{>0}$ such that $\operatorname{Re}(b_j+N+1)>0$, $1\leqslant j\leqslant p$. Then for $\ell\geqslant N+1$,
		\begin{equation}\label{pFp estimate}
			\biggl|{}_pF_p\biggl[\begin{matrix}
				a_1+\ell,\dots,a_p+\ell\\
				b_1+\ell,\dots,b_p+\ell
			\end{matrix};z\biggr]\biggr|\leqslant C\me^{\gamma|z|},
		\end{equation}
		where $\gamma>0$ is a constant independent of $\ell$, $N$ and $z$.
\end{Remark}
